\documentclass[12pt,a4paper]{article}
\usepackage[utf8]{inputenc}
\usepackage[T1]{fontenc}
\usepackage[brazil]{babel}
\usepackage{graphicx}
\usepackage{url}
\usepackage{geometry}
\usepackage{float}
\usepackage{booktabs}
\geometry{a4paper, left=3cm, right=2cm, top=3cm, bottom=2cm}

\begin{document}

\begin{center}
    \textbf{RELATÓRIO PARCIAL DE ATIVIDADES DE INICIAÇÃO CIENTÍFICA}\\
    \textbf{Programa PIBIC-Af / Cadastro ICT}
\end{center}

\section{IDENTIFICAÇÃO}

\begin{description}
    \item[Centro:] Centro de Ciências Exatas e de Tecnologia (CCET)
    \item[Departamento:] Departamento de Computação (DC)
    \item[Orientador:] Prof. Dr. Murilo Naldi
    \item[Aluno(a):] Maylon Martins de Melo
    \item[Título do Projeto:] Desenvolvimento e Melhoria do MustaCHE através do CORE-SG
    \item[Período de Vigência:] Setembro/2025 a Agosto/2026
\end{description}

\section{INTRODUÇÃO}

O objetivo central deste projeto de Iniciação Científica é aprimorar a eficiência computacional da ferramenta MustaCHE (\textit{Multiple Cluster Hierarchies Explorer}) através da integração do algoritmo CORE-SG. O software original, embora cientificamente relevante, encontrava-se em estado de obsolescência tecnológica (desenvolvido em Python 2.7 e Java), o que dificultava sua manutenção e execução em ambientes modernos.

Os objetivos específicos para este primeiro período foram:
\begin{itemize}
    \item Realizar a reengenharia completa da plataforma (\textit{MustaCHE v2});
    \item Modernizar a \textit{stack} tecnológica para Python 3.11, Flask e Plotly.js;
    \item Implementar containerização via Docker para garantir reprodutibilidade;
    \item Restaurar e validar as funcionalidades legadas (HAI Matrix, Meta-Clustering) antes da integração do CORE-SG.
\end{itemize}

\section{RESUMO DAS ATIVIDADES E RESULTADOS OBTIDOS}

Durante o período relatado, as atividades concentraram-se na reconstrução da arquitetura do sistema. A metodologia adotada focou na modularização e escalabilidade. Os principais resultados incluem:

\begin{enumerate}
    \item \textbf{Reengenharia do Backend}: Migração integral do código legado para Python 3.11. Removeu-se a dependência de componentes Java, centralizando a lógica de clustering na biblioteca \textit{scikit-learn} com o algoritmo HDBSCAN.
    \item \textbf{Interface Gráfica Interativa}: Desenvolvimento de um novo \textit{dashboard} utilizando Plotly.js, capaz de renderizar simultaneamente Dendrogramas, Matrizes de Similaridade HAI e \textit{Reachability Plots} interativos (Figura \ref{fig:dashboard}).
    \item \textbf{Containerização}: Criação de um ambiente Docker multi-estágio que automatiza a compilação de módulos Cython e a instalação de dependências científicas, garantindo que a ferramenta funcione de forma idêntica em qualquer plataforma.
    \item \textbf{Implementação da Matriz HAI}: Restauração do cálculo do \textit{Hierarchy Agreement Index} (Figura \ref{fig:hai}), essencial para comparar diferentes hierarquias de cluster e identificar a estabilidade dos parâmetros fundamentais do HDBSCAN ($m_{pts}$).
    \item \textbf{Visualização Hierárquica (Dendrograma)}: Desenvolvimento de um componente interativo que permite visualizar a estrutura de agrupamento e realizar cortes dinâmicos ("thresholds") para a definição de meta-clusters (Figura \ref{fig:dendrogram}).
\end{enumerate}

\begin{figure}[H]
    \centering
    \includegraphics[width=0.9\textwidth]{img/batch results.png}
    \caption{Interface principal do MustaCHE v2 com resultados de processamento em lote.}
    \label{fig:dashboard}
\end{figure}

\begin{figure}[H]
    \centering
    \includegraphics[width=0.65\textwidth]{img/reach plots.png}
    \caption{Visualização de Reachability Plots interativos no sistema modernizado.}
    \label{fig:reachability}
\end{figure}

\begin{figure}[H]
    \centering
    % Substitua pelo caminho da sua imagem da Matriz HAI
    % \includegraphics[width=0.8\textwidth]{img/hai_matrix.png}
    \fbox{\begin{minipage}[c][6cm]{0.8\textwidth}
        \centering \vspace{2.5cm} [Inserir aqui: Captura de tela da Matriz Similarity HAI de Hierarquias]
    \end{minipage}}
    \caption{Matriz de Similaridade HAI, permitindo a análise de estabilidade das hierarquias geradas.}
    \label{fig:hai}
\end{figure}

\begin{figure}[H]
    \centering
    % Substitua pelo caminho da sua imagem do Dendrograma
    % \includegraphics[width=0.8\textwidth]{img/dendrogram.png}
    \fbox{\begin{minipage}[c][6cm]{0.8\textwidth}
        \centering \vspace{2.5cm} [Inserir aqui: Captura de tela do Dendrograma com Linha de Corte]
    \end{minipage}}
    \caption{Dendrograma de Meta-Clusters com funcionalidade de seleção dinâmica de medoids.}
    \label{fig:dendrogram}
\end{figure}

\section{DIFERENCIAIS TÉCNICOS E CIENTÍFICOS}

Para além da reconstrução funcional, o projeto introduziu melhorias críticas no fluxo de trabalho científico da ferramenta:

\begin{enumerate}
    \item \textbf{Reprodutibilidade Científica via Docker}: O uso de containerização de múltiplos estágios automatiza a compilação de módulos legados em Cython. Isso elimina a barreira do "it works on my machine", permitindo que qualquer pesquisador reproduza as análises independentemente do sistema operacional utilizado.
    \item \textbf{MustaCHE como Biblioteca Python}: O software foi reestruturado para ser não apenas um utilitário web, mas um pacote Python instalável. Isso permite que outros cientistas de dados importem suas funções (ex: \texttt{from mustache.core import run_clustering}) em seus próprios fluxos de trabalho de pesquisa via \textit{Jupyter Notebooks} ou scripts, aumentando significativamente o impacto acadêmico da ferramenta.
    \item \textbf{Padronização e Metadados de Citação}: Implementou-se o suporte ao formato \textit{Citation File Format} (\texttt{CITATION.cff}) no repositório GitHub. Isso permite que a ferramenta seja automaticamente citable no GitHub e em serviços como o \textit{Zenodo}, garantindo que o crédito acadêmico seja atribuído corretamente aos autores originais (Neto et al., 2018).
\end{enumerate}

\section{PRODUÇÃO TÉCNICO-CIENTÍFICA}

As atividades desenvolvidas resultaram na criação de um repositório de código aberto no GitHub (\url{https://github.com/Maylon-hub/mustache}), estruturado com licença \textbf{BSD 3-Clause}. Esta licença foi escolhida por ser o padrão ouro do ecossistema científico Python (\textit{NumPy}, \textit{scikit-learn}), permitindo o uso acadêmico e comercial livre, enquanto protege os nomes dos autores originais e da UFSCar. Além do código-fonte, o repositório contém a documentação técnica necessária para reprodução integral dos experimentos.

\section{CRONOGRAMA E PRÓXIMOS PASSOS}

Para o restante do período de vigência da bolsa, o plano de trabalho foca na \textbf{Fase 2: Otimização e Comparação}. O diferencial desta etapa será a avaliação quantitativa do algoritmo CORE-SG em comparação ao HDBSCAN padrão, utilizando métricas como o índice \textbf{DBCV (Density-Based Clustering Validation)} e o tempo de execução (\textit{throughput}) em conjuntos de dados de larga escala.

\begin{table}[H]
    \centering
    \caption{Cronograma de atividades para os próximos seis meses.}
    \begin{tabular}{@{}ll@{}}
    \toprule
    \textbf{Mês} & \textbf{Atividade Planejada} \\ \midrule
    Mês 7-8 & Integração do CORE-SG e refatoração do grafo de similaridade. \\
    Mês 9-10 & Bateria de experimentos comparativos e coleta de métricas (DBCV, Tempo). \\
    Mês 11 & Análise estatística dos resultados e escrita do relatório final. \\
    Mês 12 & Submissão de resultados a eventos científicos/periódicos (ex: Congresso de IC). \\ \bottomrule
    \end{tabular}
\end{table}

\section{AUTO-AVALIAÇÃO ASSINADA}

Avalio minha participação neste primeiro período como produtiva e de grande aprendizado técnico. A superação dos desafios de obsolescência do código original exigiu o domínio de ferramentas modernas de desenvolvimento (Docker, Flask, Git) e uma compreensão profunda da lógica matemática do HDBSCAN e do índice HAI. Sinto-me tecnicamente preparado para iniciar a integração do algoritmo CORE-SG na próxima fase do projeto.

\vspace{1.5cm}
\begin{center}
    \rule{10cm}{0.5pt}\\
    \textbf{Maylon Martins de Melo}\\
    Aluno(a)
\end{center}

\section{AVALIAÇÃO DO(A) ORIENTADOR(A) ASSINADA}

O aluno Maylon demonstrou excelente desempenho, proatividade e capacidade de resolução de problemas complexos de software. A entrega de uma versão modernizada da ferramenta MustaCHE superou as expectativas iniciais, fornecendo uma base sólida e citable para as próximas etapas da pesquisa. O aluno cumpriu rigorosamente o cronograma de atividades proposto.

\vspace{1.5cm}
\begin{center}
    \rule{10cm}{0.5pt}\\
    \textbf{Prof. Dr. Murilo Naldi}\\
    Orientador(a)
\end{center}

\end{document}
